% Generated draft: requires mathematical and source review.
\documentclass[11pt]{article}
\usepackage[margin=25mm]{geometry}
\usepackage{fontspec}
\setmainfont{Latin Modern Roman}
\usepackage{amsmath,amssymb,mathtools,unicode-math}
\setmathfont{Latin Modern Math}
\usepackage{enumitem,needspace,xurl,hyperref}
\setlength{\parindent}{0pt}
\setlength{\parskip}{5pt}
\setlength{\emergencystretch}{3em}
\newcommand{\N}{\mathbb{N}}
\newcommand{\Z}{\mathbb{Z}}
\newcommand{\Q}{\mathbb{Q}}
\newcommand{\R}{\mathbb{R}}
\newcommand{\C}{\mathbb{C}}
\newcommand{\F}{\mathbb{F}}
\newcommand{\A}{\mathbb{A}}
\newcommand{\PP}{\mathbb{P}}
\newcommand{\E}{\mathbb{E}}
\newcommand{\eps}{\varepsilon}
\newcommand{\dd}{\,\mathrm{d}}
\begin{document}
\section*{UnsolvedMath ID: 701398}
\textbf{Source problem code:} catalogue entry
\subsection{Definitions and mathematical statement}
A variety of groups is a class consisting of all groups satisfying a specified set of group identities. A group identity is an equation between group words that holds for every assignment of its variables to elements of the group. A subvariety of a variety \(\mathcal V\) is a variety of groups contained in \(\mathcal V\). Write
\[
\operatorname{Sub}(\mathcal V)=\{\mathcal W: \mathcal W\text{ is a variety of groups and }\mathcal W\subseteq\mathcal V\}.
\]
Does there exist a variety of groups \(\mathcal V\) such that
\[
\lvert\operatorname{Sub}(\mathcal V)\rvert=\aleph_0
\]
and there is a sequence \((\mathcal V_n)_{n\geq 0}\) of subvarieties of \(\mathcal V\) satisfying
\[
\mathcal V_{n+1}\subsetneq\mathcal V_n\qquad\text{for every integer }n\geq 0?
\]
Here proper containment is containment of classes of groups. The required chain already forces infinitely many distinct subvarieties, so interpreting ``only countably many'' as either at most countably many or countably infinitely many gives the same existence question.
\subsection{Short English statement}
Does some variety of groups have exactly countably infinitely many subvarieties while admitting an infinite strictly descending sequence of subvarieties?
\subsection{Sources}
\begin{itemize}
\item Supplied local UnsolvedMath registry record, permanent ID \(701398\), source code \texttt{catalogue entry}, titled ``Kourovka Notebook Problem 17.96.''
\item The supplied record attributes Problem 17.96 to P. M. Neumann, proposed in 2010, in \emph{The Kourovka Notebook (Unsolved Problems in Group Theory)}, 17th issue (2010), with reference also to the 21st issue edition (2026), arXiv:1401.0300v48. The collection description names E. I. Khukhro and V. D. Mazurov as editors.
\item The registry labels the problem ``open.'' This is a supplied status label, not an independently verified assessment of the literature. The cited source has not been independently consulted for this draft.
\end{itemize}
\end{document}
